% Generated deterministically from the audited Markdown source.
% Responsible author: Carptopus.
\documentclass[11pt]{article}
\usepackage[a4paper,margin=27mm]{geometry}
\usepackage[T1]{fontenc}
\usepackage[utf8]{inputenc}
\usepackage{lmodern}
\usepackage{microtype}
\usepackage{amsmath,amssymb,amsthm,mathtools}
\usepackage{needspace}
\usepackage{xcolor}
\usepackage[colorlinks=true,linkcolor=blue!55!black,citecolor=blue!55!black,urlcolor=blue!65!black]{hyperref}
\usepackage{xurl}
\usepackage{fancyhdr}

\newtheorem{theorem}{Theorem}
\newtheorem{corollary}[theorem]{Corollary}

\pagestyle{fancy}
\fancyhf{}
\fancyhead[L]{Carptopus}
\fancyhead[R]{A Walsh-convolution proof}
\fancyfoot[C]{\thepage}
\setlength{\headheight}{14pt}
\setlength{\parindent}{0pt}
\setlength{\parskip}{0.44em}
\setlength{\emergencystretch}{4em}
\urlstyle{same}

\title{A Walsh-convolution proof of two conjectures of Lachaume}
\author{Carptopus\\\href{mailto:carptopus@163.com}{carptopus@163.com}}
\date{4 September 2026}

\hypersetup{
  pdftitle={A Walsh-convolution proof of two conjectures of Lachaume},
  pdfauthor={Carptopus},
  pdfsubject={Lachaume's derivative convolution as a Walsh symmetric additive convolution},
  pdfkeywords={Walsh convolution, symmetric additive convolution, real-rootedness, root location, hyperbolic polynomials}
}

\begin{document}
\maketitle

\begin{center}
\fcolorbox{black!35}{black!4}{\parbox{0.88\textwidth}{\centering\small
Public Beta manuscript, version 0.1-beta. The mathematical proof and complete
manuscript have passed project-internal zero-trust audits. External mathematical review
and formal peer review are pending.}}
\end{center}

\begin{abstract}


For degree-\(p\) polynomials \(P\) and \(Q\), consider

\[
\mathcal D_p(P,Q)(x)=\sum_{k=0}^{p}P^{(k)}(x)Q^{(p-k)}(x).
\]

Lachaume conjectured that this polynomial is real-rooted when \(P\) and \(Q\)
are real-rooted, and more generally that every zero of \(\mathcal D_p(P,Q)\)
lies in the convex hull of the zeros of \(P\) and \(Q\). We observe the exact
identity

\[
\mathcal D_p(P,Q)(x)=p!\,(P\boxplus_p Q)(2x),
\]

where \(\boxplus_p\) is the classical Walsh symmetric additive convolution.
Walsh's root-location theorem then proves both assertions in every degree.
Through Lachaume's equivalences, this also establishes the corresponding
semi-symmetric hyperbolicity statement. The convolution and its root-location
theorem are classical; the purpose of this note is to record their direct
application to Lachaume's questions.

\end{abstract}

\section{The derivative sum}


For complex polynomials \(P,Q\) of degree at most \(p\), where \(p\ge 1\),
define

\[
\mathcal D_p(P,Q)(x)
=\sum_{k=0}^{p}P^{(k)}(x)Q^{(p-k)}(x).
\]

The degree-\(p\) Walsh symmetric additive convolution, in the normalization of
Marcus, Spielman and Srivastava \cite{mss}, is

\[
(P\boxplus_p Q)(y)
=\frac1{p!}\sum_{k=0}^{p}P^{(k)}(y)Q^{(p-k)}(0).
\]

\Needspace{0.16\textheight}
\begin{theorem}
\label{thm:identity}

For all complex polynomials \(P,Q\) of degree at most \(p\),

\[
\mathcal D_p(P,Q)(x)=p!\,(P\boxplus_pQ)(2x).
\]

\end{theorem}

\begin{proof}

Fix \(x\) and set

\[
F(t)=\sum_{k=0}^{p}P^{(k)}(x+t)Q^{(p-k)}(x-t).
\]

Differentiating gives

\[
\begin{aligned}
F'(t)
&=\sum_{k=0}^{p}P^{(k+1)}(x+t)Q^{(p-k)}(x-t)\\
&\quad-\sum_{k=0}^{p}P^{(k)}(x+t)Q^{(p-k+1)}(x-t)=0.
\end{aligned}
\]

The interior terms cancel after shifting the index, while the two endpoint
terms vanish because the \((p+1)\)-st derivatives are zero. Hence \(F\) is
constant. Evaluating at \(t=0\) and \(t=x\) yields

\[
\mathcal D_p(P,Q)(x)
=\sum_{k=0}^{p}P^{(k)}(2x)Q^{(p-k)}(0)
=p!\,(P\boxplus_pQ)(2x).
\]

\end{proof}

\section{Root location}


The classical Walsh theorem says that if the zeros of \(P\) and \(Q\) lie in
closed discs \(C_1\) and \(C_2\), respectively, then every zero of
\(P\boxplus_pQ\) lies in the Minkowski sum \(C_1+C_2\). Leake and Ryder
\cite[p.~1]{leake-ryder} state this complex-disc form explicitly and trace the convolution
to Walsh \cite{walsh}.

\Needspace{0.16\textheight}
\begin{theorem}
\label{thm:root-location}

Let \(P,Q\in\mathbb C[X]\) have degree \(p\). Every zero of
\(\mathcal D_p(P,Q)\) lies in the convex hull of all zeros of \(P\) and \(Q\).

\end{theorem}

\begin{proof}

Let \(S\) be the union of the two finite zero multisets, and let \(C\) be any
closed disc containing \(S\). Walsh's theorem puts every zero of
\(P\boxplus_pQ\) in \(2C=C+C\). Theorem~\ref{thm:identity} therefore puts every zero of
\(\mathcal D_p(P,Q)\) in \(C\).

The intersection of all closed discs containing a finite planar set \(S\) is
its convex hull. One inclusion follows from convexity. For the other, strictly
separate any point outside \(\operatorname{conv}(S)\) by a line, then
approximate the half-plane containing \(S\) by a sufficiently large disc whose
center recedes along the opposite normal. Thus every zero of
\(\mathcal D_p(P,Q)\) belongs to \(\operatorname{conv}(S)\).
\end{proof}


\Needspace{0.16\textheight}
\begin{corollary}
\label{cor:real-rooted}

If \(P\) and \(Q\) are real-rooted, then \(\mathcal D_p(P,Q)\) is real-rooted.

\end{corollary}

\begin{proof}

All input zeros lie on the real line, so their convex hull is a real interval.
\end{proof}


\section{Lachaume's conjectures}


Lachaume \cite[Theorem~4.9]{lachaume} states that his Conjecture 4 is equivalent to
Corollary~\ref{cor:real-rooted} for every degree \(p\). Therefore Conjecture 4 holds for all
\(p\ge1\). By \cite[Theorem~4.8]{lachaume}, this proves the corresponding hyperbolicity of
the semi-symmetric polynomials \(s_{n,p}\) for every admissible \(n\ge p\).

Lachaume's Conjecture 5 is precisely Theorem~\ref{thm:root-location} and hence also holds for every
\(p\ge1\). The paper further records that Conjecture 4 implies Conjecture 2 in
its stated nonnegative-coefficient setting. This concavity implication is not
new: Lachaume's later Lovelock-gravity paper explicitly records the
real-rooted diagonal criterion as an application of Walsh's theorem
\cite[Theorem~2.7(d)]{lachaume-lovelock}. If Lachaume's diagonal
restriction polynomial \(\bar f_{\vec a}\) has actual degree \(q<p\) because
its leading coefficients vanish, the degree-\(q\) result gives concavity of
\(f^{1/q}\); composing with the increasing concave map
\(u\mapsto u^{q/p}\) gives concavity of \(f^{1/p}\). The positive constant
case is immediate. The identically zero function lies outside Lachaume's
strict positive-value definition of \(\mu\)-concavity; it is concave only
under the usual extended nonnegative-function convention.

These implications do not address Lachaume's Conjectures 1 or 3.

\section{Scope and prior work}


Walsh introduced the symmetric additive convolution and proved its classical
root-location property \cite{walsh}. Marcus, Spielman and Srivastava \cite[Definition~1.1 and Theorem~1.3]{mss} use the same convolution with a \(1/p!\) normalization and
record its real-rootedness preservation. Leake and Ryder \cite{leake-ryder} use the
unnormalized convention and explicitly recall the Minkowski-sum statement for
complex discs. Multiplication by \(p!\) does not affect zeros.

Lachaume's later paper \cite[Theorem~2.7(d)]{lachaume-lovelock} already invokes Walsh's theorem to
deduce concavity from real-rootedness of a diagonal restriction. Thus the
broad Lachaume--Walsh connection is prior work. Theorem~\ref{thm:identity} is an elementary
expression of translation covariance, not a new convolution theorem. The
contribution of this note is limited to the explicit identity with its exact
scaling, its application to Lachaume's Conjectures 4 and 5, and the complex
convex-hull deduction. Searches completed on 4 September 2026 found no earlier
source stating this identity or using it to settle both conjectures. That
search result does not establish global priority, and no claim of priority for
Walsh's theorem is made.

\section*{Computational check}


The accompanying symbolic script verifies Theorem~\ref{thm:identity}, the telescoping
certificate, and the apolar coefficient identity for generic polynomials of
degrees \(1\) through \(8\). It also contains a destructive control obtained by
deleting one endpoint term. The script is supplementary evidence; the proof
above is independent of it.

\section*{AI-assisted research disclosure}


Carptopus is the responsible author and takes responsibility for the claims,
proofs, source use, and released artifacts. OpenAI Codex was used for
AI-assisted literature organization, proof development and stress testing,
exact verification, adversarial auditing, and manuscript preparation.

\begin{thebibliography}{9}
\footnotesize
\raggedright

\bibitem{lachaume}
X.~Lachaume,
\emph{On the concavity of a sum of elementary symmetric polynomials},
arXiv:1712.10327v1 (2017).
\url{https://arxiv.org/abs/1712.10327}.

\bibitem{mss}
A.~W. Marcus, D.~A. Spielman and N.~Srivastava,
\emph{Finite free convolutions of polynomials},
Probability Theory and Related Fields \textbf{182} (2022), 807--848.
\href{https://doi.org/10.1007/s00440-021-01058-2}{doi:10.1007/s00440-021-01058-2}.

\bibitem{leake-ryder}
J.~Leake and N.~Ryder,
\emph{Connecting the q-multiplicative convolution and the finite difference convolution},
Advances in Mathematics \textbf{374} (2020), 107332.
\href{https://doi.org/10.1016/j.aim.2020.107332}{doi:10.1016/j.aim.2020.107332}.

\bibitem{walsh}
J.~L. Walsh,
\emph{On the location of the roots of certain types of polynomials},
Transactions of the American Mathematical Society \textbf{24} (1922), 163--180.

\bibitem{lachaume-lovelock}
X.~Lachaume,
\emph{The constraint equations of Lovelock gravity theories: a new
$\sigma_k$-Yamabe problem},
Journal of Mathematical Physics \textbf{59} (2018), 072501.
\href{https://doi.org/10.1063/1.5023758}{doi:10.1063/1.5023758}.

\end{thebibliography}

\end{document}
